\chapter{Creio em Jesus Cristo, seu Filho único, nosso Senhor}
\chaptermark{Creio em Jesus Cristo}

Este capítulo aprofunda o segundo artigo do Símbolo Apostólico, professando a fé em Jesus Cristo como Filho único de Deus Pai, nosso Senhor. Desdobramos os parágrafos correspondentes do Catecismo da Igreja Católica (CIC), explorando sua identidade divina, títulos reveladores, concepção virginal, vida oculta e início da vida pública. Para adultos em processo de iniciação cristã pelo RICA, essa revelação de Cristo desperta o desejo de segui-lo, preparando o coração para os sacramentos que nos unem a Ele na Páscoa.

\section{Jesus Cristo, Filho Unigênito de Deus (\S 422--455)}

A fé cristã centra-se em Jesus Cristo, verdadeiro Deus e verdadeiro homem, Filho eterno do Pai. O Novo Testamento testemunha unanimemente sua filiação divina, desde o anúncio do anjo a Maria até a confissão de Tomé: ``Meu Senhor e meu Deus!''. Essa verdade é o cerne do Evangelho, professado no Credo como artigo fundamental.

Jesus é o Filho unigênito, gerado não criado, consubstancial ao Pai pela eternidade. Os Evangelhos Sinóticos e de João revelam-no como Messias prometido, cumprindo as Escrituras hebraicas, enviado para salvar Israel e todas as nações.

A Igreja primitiva, guiada pelo Espírito, formula essa fé nos Concílios: Niceia (325) afirma ``Deus de Deus, Luz de Luz, Deus verdadeiro de Deus verdadeiro'', combatendo heresias como o arianismo, que negava sua divindade plena.

Os Apóstolos e Pais da Igreja transmitem essa querigma viva, convidando à adesão pessoal: crer em Cristo é participar de sua vida filial.

Para adultos a partir dos 20 anos no RICA, essa seção inicia a catequese cristológica, convidando à profissão de fé em Jesus como Senhor no rito de aceitação, fomentando a conversão que reconhece Nele o Salvador pessoal.

\section{Os Títulos e Nomes de Jesus Cristo (\S 456--483)}

Jesus recebe múltiplos títulos que revelam sua identidade: Cristo (Messias, Ungido), Filho de Davi, Filho do Homem, Filho de Deus. Cada nome ilumina sua missão messiânica, real, profética e sacerdotal, cumprida na Cruz e Ressurreição.

``Filho de Deus'' expressa a relação íntima e eterna com o Pai, manifestada nos milagres, exorcismos e perdão dos pecados. Os demônios o confessam primeiro, forçados por sua autoridade divina.

``Senhor'' (Kyrios) evoca a divindade de YHWH no Antigo Testamento, aplicado a Jesus como ressuscitado e Pantocrator. Paulo o proclama Senhor de todos, a quem todo joelho se dobrará.

Esses nomes não são meros apelidos, mas confissões de fé que a Igreja adora na liturgia, unindo céu e terra em louvor.

Na catequese para adultos não iniciados, meditar nesses títulos enriquece as celebrações da Palavra no catecumenato, ajudando a invocar Jesus como Senhor na oração diária e preparar o testemunho público no rito de eleição.

\section{A Concepção Virginal e o Nascimento de Jesus (\S 484--511)}

A concepção de Jesus é obra do Espírito Santo no seio virginal de Maria: ``O Espírito Santo descerá sobre ti e o poder do Altíssimo te cobrirá com sua sombra.'' Maria, cheia de graça, consente livremente, tornando-se Mãe de Deus (Theotokos).

Esse mistério afirma a humanidade plena de Cristo, assumida desde a concepção, sem pecado original. O nascimento em Belém, humilde e glorioso, cumpre as profecias: ``Uma virgem conceberá e dará à luz um filho''.
A Encarnação é o ápice da história da salvação: o Verbo se fez carne para nos salvar, unindo divindade e humanidade em uma só Pessoa hipostática, definida em Calcedônia (451).

A Virgem Maria, nova Eva, coopera na Redenção, modelo de fé para a Igreja nascente.

Para catecúmenos adultos, contemplar a Encarnação no tempo de purificação quaresmal inspira confiança na ação do Espírito, preparando para o Batismo que nos torna filhos no Filho único.

\section{A Vida Oculta e as Tentações de Jesus (\S 512--534)}

Após o nascimento, Jesus vive 30 anos oculto em Nazaré, em obediência filial a José e Maria, revelando a santidade na vida cotidiana. Essa fase humilde ensina que a grandeza divina se manifesta no ordinário.

O Batismo no Jordão marca o início público: os céus se abrem, o Espírito desce como pomba, o Pai proclama: ``Tu és o meu Filho amado''. João Batista o aponta como Cordeiro de Deus.

Nas tentações no deserto, Jesus, cheio do Espírito, rejeita as seduções de Satanás com a Palavra de Deus, recapitulado Adão vitorioso, iniciando sua missão de libertação.

Esses eventos inauguram o tempo messiânico, convidando ao arrependimento e à fé no Reino.

Em catequese RICA adulta, a vida oculta de Jesus modela a paciência no catecumenato longo, enquanto as tentações armam contra provações, com exorcismos menores fortalecendo para a Vigília Pascal.

\section{O Início da Vida Pública de Jesus (\S 535--570)}

Jesus inicia sua pregação na Galileia, chamando discípulos e anunciando o Reino de Deus: ``O tempo se cumpriu e o Reino de Deus está próximo; convertei-vos e crede no Evangelho''. Batiza, perdoa pecados e realiza milagres como sinais do Reino.

Os milagres --- curas, multiplicação dos pães, domínio sobre a natureza --- manifestam sua messianidade, culminando na Transfiguração, visão da glória divina ante Pedro, Tiago e João.

Jesus ensina em parábolas, revelando mistérios do Reino aos simples, enquanto confronta fariseus pela hipocrisia, purificando o judaísmo para a universalidade.

Sua mesa aberta a pecadores e pobres escandaliza, mas revela a misericórdia paterna, prefigurando a Eucaristia.

Para adultos em busca de iniciação, o anúncio de Jesus no período de iluminação quaresmal desperta fome pelo Reino, integrando parábolas na catequese para viver a conversão concreta rumo aos sacramentos pascais.

\section{Conclusão}

Professar fé em Jesus Cristo, Filho único e Senhor, revela o coração da salvação: Deus conosco. Adultos no RICA, sigam-No para a iniciação plena na Páscoa.